\section{Capabilities}
\label{sec:capabilities}

This section states what the system can be asked to do. \Cref{sec:fulfilment}
maps these onto the requirement identifiers of ETA-TR-2026-001, and
\Cref{sec:results} gives the evidence.

\subsection{It selects its own objectives}

No test objective enters the loop from outside. The planner reads the graph,
decides what is worth testing next, and writes an objective together with the
requirements it claims to cover. Proposals are rejected structurally and as
duplicates, using an embedding-cosine check behind a lexical pre-filter, so a
campaign does not spend its budget re-testing what it has already tried.

Requirement citation is enforced at both construction paths. An earlier defect
allowed decorated identifiers to bypass the check, which is the kind of failure
that quietly destroys a coverage claim, and it is now closed at both sites
rather than at the one where it was found.

\subsection{It runs on two platforms without changing}

One planner and one graph drive both Android and the web. The Android player
uses the accessibility service with a companion on-device application; the web
player uses an accessibility-annotated browser snapshot. The planner does not
know which it is driving.

Adding a platform meant writing a player against the same contract. Adding a
\emph{target} on an existing platform means writing a configuration profile and
nothing else: no code, no build, no redeploy.

\subsection{It works with whatever knowledge exists}

Six retrieval sources are supported and all six are optional: specification,
screen designs, flow designs, the live observed interface, defect history, and
the navigation tree. A source that is absent contributes nothing rather than
blocking the round.

The practical range this covers runs from a target with a full specification
and a defect export down to a target with nothing at all, where the system
works purely from what it observes. \Cref{sec:limits} is candid about how much
quality is lost at that end of the range.

\subsection{It learns across runs}

Four kinds of knowledge accumulate and are read back into the next round.

\begin{description}[leftmargin=2.4em,style=nextline,itemsep=3pt,topsep=4pt]
  \item[Navigation memory.] Paths that reached a screen are recorded, along
    with transitions that led nowhere, so a later test walks a known route
    instead of rediscovering it. This is the mechanism behind the fall in
    median device steps from 50 to 9.
  \item[Execution history and error patterns.] Failures are mined for recurring
    patterns, which feed a strategy memory the planner consults when choosing
    how to approach a screen.
  \item[The application model.] Interface states and the transitions between
    them are built entirely from observation, with no design export required.
  \item[Findings.] Deduplicated claims about the application, each attributed
    and each linked to the runs that support it.
\end{description}

\subsection{It reports in a form a stakeholder can act on}

A campaign is configured, started and watched from a dashboard, and exports a
report whose numbers are computed from the graph while the language model
supplies only connecting prose. Agent-caused failures are filtered out before a
defect is named.

\begin{figure}[H]
  \centering
  \begin{subfigure}{0.49\textwidth}
    \centering\includegraphics[width=\textwidth]{dashboard-live-logs.jpg}
    \caption{Live campaign view.}
  \end{subfigure}\hfill
  \begin{subfigure}{0.49\textwidth}
    \centering\includegraphics[width=\textwidth]{dashboard-app-model.jpg}
    \caption{The application model as it is built.}
  \end{subfigure}
  \caption[Operator surface]{The operator surface. These captures are from a
  development run and are included to show what the interface presents; the
  figures this report stands behind are in \Cref{sec:results}.}
  \label{fig:dashboard}
\end{figure}

\subsection{It bounds what it is allowed to do}

Operations that cannot be undone are blocked at the configuration layer rather
than discouraged in a prompt. On the production target this covers account
registration, one-time-password verification, payment and account deletion.
Across 123 recorded executions no irreversible action occurred.

Test credentials reach the player, which has to type them, and never the
planner, which receives the account role only. A secret that never enters the
planner's context cannot appear in a planner trace, a prompt log or the
dashboard.

Application content is labelled as untrusted data in every prompt that carries
it. The rule is implemented separately in the player, the investigator and the
evidence reviewer, because each assembles its own prompt and a single shared
helper would have been bypassed by whichever path did not use it.
